% Part: sets-functions-relations
% Chapter: sets
% Section: unions-and-intersections

\documentclass[../../../include/open-logic-section]{subfiles}

\begin{document}

\olfileid[fa-IR]{sfr}{set}{uni}
\olsection{اجتماع‌ها و اشتراک‌ها}

\begin{explain}
در \olref[sfr][set][bas]{sec}، تعریف مجموعه‌ها به روش تجرید، یعنی
تعریف‌هایی به‌صورت $\Setabs{x}{\phi(x)}$، را معرفی کردیم. در این روش
ویژگی~$\phi$ را به کار می‌بریم و این ویژگی می‌تواند به مجموعه‌هایی که
پیش‌تر تعریف کرده‌ایم اشاره کند. بنابراین، برای نمونه، اگر $A$ و~$B$
مجموعه باشند، مجموعهٔ $\Setabs{x}{x \in A \lor x \in B}$ از همهٔ
اشیایی تشکیل شده است که در شمار !!{element}s $A$ یا~$B$ هستند؛ یعنی
مجموعه‌ای است که !!{element}s $A$ و~$B$ را با هم گرد می‌آورد.
می‌توان این وضعیت را مانند \olref{fig:union} مجسم کرد؛ در آنجا ناحیهٔ
برجسته، !!{element}s دو مجموعهٔ $A$ و~$B$ را در کنار هم نشان می‌دهد.

\begin{figure}
  \olasset{assets/diagrams/union.tikz}
  \caption{اجتماع $A \cup B$ دو مجموعه، مجموعهٔ !!{element}s
   $A$ همراه با اعضای~$B$ است.}
  \ollabel{fig:union}
\end{figure}

این عمل روی مجموعه‌ها---یعنی باهم‌ترکیب‌کردن آنها---بسیار سودمند و
پرکاربرد است؛ ازاین‌رو نام و نمادی صوری به آن می‌دهیم.
\end{explain}

\begin{defn}[اجتماع]
\emph{اجتماع} دو مجموعهٔ $A$ و $B$ که آن را $A \cup B$ می‌نویسیم،
مجموعهٔ همهٔ اشیایی است که در شمار !!{element}s $A$ یا $B$ یا هر دو هستند.
\[
A \cup B = \Setabs{x}{x \in A \lor x \in B}
\]
\end{defn}

\begin{ex}
ازآنجاکه تعداد تکرار !!{element}s اهمیتی ندارد، اجتماع دو مجموعه که
!!a{element} مشترک دارند، آن !!{element} را تنها یک بار در بر می‌گیرد؛
برای نمونه، $\{ a, b, c\} \cup \{ a, 0, 1\} = \{a, b, c, 0, 1\}$.

اجتماع یک مجموعه با یکی از زیرمجموعه‌هایش همان مجموعهٔ بزرگ‌تر است:
$\{a, b, c \} \cup \{a \} = \{a, b, c\}$.

اجتماع یک مجموعه با مجموعهٔ تهی همان مجموعه است: $\{a,
b, c \} \cup \emptyset = \{a, b, c \}$.
\end{ex}

\begin{prob}
ثابت کنید اگر $A \subseteq B$، آنگاه $A \cup B = B$.
\end{prob}

\begin{explain}
می‌توانیم عملی ``دوگان'' با اجتماع را نیز در نظر بگیریم. این عمل
مجموعهٔ همهٔ !!{element}s را تشکیل می‌دهد که هم در شمار !!{element}s~$A$
و هم در شمار !!{element}s~$B$ هستند. این عمل \emph{اشتراک} نام دارد و می‌توان
آن را مانند \olref{fig:intersection} نمایش داد.
\begin{figure}
  \olasset{assets/diagrams/intersection.tikz}
  \caption{اشتراک $A \cap B$ دو مجموعه، مجموعهٔ
    !!{element}s مشترک آنهاست.}
  \ollabel{fig:intersection}
\end{figure}
\end{explain}

\begin{defn}[اشتراک]
\emph{اشتراک} دو مجموعهٔ $A$ و $B$ که آن را $A \cap B$ می‌نویسیم،
مجموعهٔ همهٔ اشیایی است که در شمار !!{element}s هر دو مجموعهٔ $A$ و~$B$ هستند.
\[
A \cap B = \Setabs{x}{x \in A \land x \in B}
\]
دو مجموعه را \emph{جدا از هم} می‌نامیم هرگاه اشتراکشان تهی باشد؛ یعنی
هیچ !!{element}s مشترکی نداشته باشند.
\end{defn}

\begin{ex}
اگر دو مجموعه هیچ !!{element}s مشترکی نداشته باشند، اشتراکشان تهی است:
$\{ a, b, c\} \cap \{ 0, 1\} = \emptyset$.

اگر دو مجموعه !!{element}s مشترک داشته باشند، اشتراکشان مجموعهٔ همهٔ
آن اعضاست: $\{a, b, c \} \cap \{a, b, d \} = \{a, b\}$.

اشتراک یک مجموعه با یکی از زیرمجموعه‌هایش همان مجموعهٔ کوچک‌تر است:
$\{a, b, c\} \cap \{a, b\} = \{a, b\}$.

اشتراک هر مجموعه با مجموعهٔ تهی، تهی است: $\{a, b, c \}
\cap \emptyset = \emptyset$.
\end{ex}

\begin{prob}
به‌دقت ثابت کنید اگر $A \subseteq B$، آنگاه $A \cap B = A$.
\end{prob}

\begin{explain}
می‌توانیم اجتماع یا اشتراک بیش از دو مجموعه را نیز تشکیل دهیم. راهی
ظریف برای بررسی حالت کلی چنین است: فرض کنید همهٔ مجموعه‌هایی را که
می‌خواهید اجتماع (یا اشتراک) آنها را تشکیل دهید، در یک مجموعه گرد
آورید. سپس می‌توانیم اجتماع همهٔ مجموعه‌های نخستین را مجموعهٔ همهٔ
اشیایی تعریف کنیم که دست‌کم به یک !!{element} از آن مجموعه تعلق دارند،
و اشتراک را مجموعهٔ همهٔ اشیایی که به هر !!{element} آن مجموعه تعلق دارند.
\end{explain}

\begin{defn}
اگر $A$ مجموعه‌ای از مجموعه‌ها باشد، آنگاه $\bigcup A$ مجموعهٔ همهٔ
اشیایی است که در شمار !!{element}sِ !!{element}s~$A$ هستند:
\begin{align*}
\bigcup A & = \Setabs{x}{x \text{ متعلق به !!a{element}ی از } A},
\text{ یعنی}\\
& = \Setabs{x}{\text{یک } B \in A
  \text{ وجود دارد به‌طوری‌که } x \in B}
\end{align*}
\end{defn}

\begin{defn}
اگر $A$ مجموعه‌ای از مجموعه‌ها باشد، آنگاه $\bigcap A$ مجموعهٔ اشیایی
است که همهٔ اعضای~$A$ در آنها مشترک‌اند:
\begin{align*}
\bigcap A & = \Setabs{x}{x \text{ متعلق به هر !!{element} از } A},
\text{ یعنی}\\
 & = \Setabs{x}{\text{برای هر } B \in A, x \in B}
\end{align*}
\end{defn}

\begin{ex}
فرض کنید $A = \{ \{ a, b \}, \{ a, d, e \}, \{ a, d \} \}$.
آنگاه $\bigcup A = \{ a, b, d, e \}$ و $\bigcap A = \{ a \}$.
\end{ex}
\begin{prob}
	نشان دهید اگر $A$ یک مجموعه و $A \in B$ باشد، آنگاه $A \subseteq \bigcup B$.
\end{prob}

می‌توان همین کار را برای دنباله‌ای از مجموعه‌ها مانند $A_1$، $A_2$، \dots
نیز انجام داد:
\begin{align*}
\bigcup_i A_i & = \Setabs{x}{x \text{ متعلق به یکی از مجموعه‌های } A_i}\\
\bigcap_i A_i & = \Setabs{x}{x \text{ متعلق به همهٔ مجموعه‌های } A_i}.
\end{align*}

هنگامی که یک \emph{مجموعهٔ اندیس} مانند $I$ داریم و مجموعه‌های مورد
نظر را با $A_i$ نشان می‌دهیم، به‌گونه‌ای که برای هر $i \in I$ یک
مجموعه در نظر گرفته شده است، می‌توانیم از صورت‌های اختصاری زیر نیز
استفاده کنیم:
\begin{align*}
	\bigcup_{i \in I} A_i & = \bigcup \Setabs{A_i }{i \in I}\\
	\bigcap_{i \in I} A_i & = \bigcap\Setabs{A_i}{i \in I}
\end{align*}

سرانجام، ممکن است بخواهیم مجموعهٔ همهٔ !!{element}s~$A$ را در نظر
بگیریم که در~$B$ نیستند. می‌توان آن را مانند \olref{difference} نمایش داد.

\begin{figure}
  \olasset{assets/diagrams/difference.tikz}
  \caption{تفاضل $A \setminus B$ دو مجموعه، مجموعهٔ
    !!{element}s~$A$ است که !!{element}s~$B$ نیستند.}
  \ollabel{difference}
\end{figure}

\begin{defn}[تفاضل]
\emph{تفاضل مجموعه‌ای}~$A \setminus B$ مجموعهٔ همهٔ !!{element}s
$A$ است که در شمار !!{element}s~$B$ نیستند؛ یعنی
\[
A\setminus B = \Setabs{x}{x\in A \text{ و } x \notin B}.
\]
\end{defn}

\begin{prob}
	ثابت کنید اگر $A \subsetneq B$، آنگاه $B \setminus A \neq \emptyset$.
\end{prob}

\end{document}
